\chapter{Reading a paper, and two projects}\label{ch:meeting}

The edition ends with work. This chapter keeps one method of reading, practised on a theorem proved in this volume, and two projects. Both projects use tools the volume has proved, and both stop on a line that is named and open.

\section{Three passes}

A research paper is not a textbook chapter. Its first page is written for a reader who must decide whether the result is the one they need; its middle is written for the reader who will check a lemma; its last page says which tradition the argument comes from. Reading it in one pass, from the first line to the last, is the wrong tool.

\emph{The first pass} asks what is claimed. Read the statement of the main theorem and the shape of the hypotheses, and not the proof. For Maier's theorem, Theorem~\ref{thm:maier}, the claim is that for every $\lambda>1$ the limsup of $(\pi(x+(\log x)^\lambda)-\pi(x))/(\log x)^{\lambda-1}$ is strictly above $1$ and the liminf is strictly below $1$. The hypotheses are none: the theorem is unconditional. The input it does use, Gallagher's prime number theorem for one large modulus, has its zero-free hypothesis in its own heading.

\emph{The second pass} asks what the skeleton is. Maier's proof builds a matrix whose rows are short intervals and whose columns are residue classes modulo a primorial. Column by column the count is Gallagher's theorem. The rows on which the Buchstab weights are large, and the rows on which they are small, give the limsup and the liminf. A reader who can say that sentence can say what the proof is for.

\emph{The third pass} checks a lemma. The lemma worth checking is the deletion of one prime factor so that the modulus has no exceptional zero, Lemma~\ref{lem:goodmoduli}. It uses the repulsion of exceptional zeros from Volume~14, and it uses nothing else.

The same three passes apply to Selberg's central limit theorem, Theorem~\ref{thm:selberg}. The claim is a Gaussian with variance $\frac12\log\log T$, at a random height, with no Riemann hypothesis in the heading. The skeleton is: move slightly off the line, replace the logarithm by a short prime sum, and compare the two in mean square. The lemma worth checking is the mean-square comparison, which is where a zero-density estimate would have entered if the argument had needed one. It does not.

\section{Project: prime factors in a short interval}

For $y\ge2$ let $\omega_y(n)$ be the number of primes $p\le y$ dividing $n$, and let $\mu_y=\sum_{p\le y}1/p$.

\begin{proposition}\label{prop:tkshort}
For real $x\ge0$, $h\ge1$ and $y\ge2$,
\[
\sum_{x<n\le x+h}\bigl(\omega_y(n)-\mu_y\bigr)^2\le h\,\mu_y+5\,\pi(y)^2.
\]
In particular, if $y\le\sqrt h$, the left side is at most $h(\mu_y+5)$.
\end{proposition}
\begin{proof}
Let $N_d$ be the number of multiples of $d$ in $(x,x+h]$, so $N_d=h/d+\theta_d$ with $|\theta_d|\le1$. Then $\sum_n\omega_y(n)=\sum_p N_p$ and $\sum_n\omega_y(n)^2=\sum_p N_p+\sum_{p\ne q}N_{pq}$, the sums over primes $p,q\le y$. Expanding,
\[
\sum_n(\omega_y-\mu_y)^2=\sum_p N_p+\sum_{p\ne q}N_{pq}-2\mu_y\sum_p N_p+\mu_y^2 N_1.
\]
The main terms are at most $h\mu_y$, because $\sum_{p\ne q}1/(pq)\le\mu_y^2$. The errors are at most $\pi(y)+\pi(y)^2+2\mu_y\pi(y)+\mu_y^2\le5\pi(y)^2$.
\end{proof}

With $h=x^\theta$ and $y=x^{\theta/2}$, an integer near $x$ has $O_\theta(1)$ prime factors larger than $y$, and $\mu_y=\log\log x+O_\theta(1)$. The proposition gives the normal order of $\omega(n)$ in intervals of length $x^\theta$. That is the proved step.

The open line is the short-interval Erd\H os--Kac theorem: for which $h=h(x)$, in particular for $h=(\log x)^A$, does $(\omega(n)-\log\log x)/\sqrt{\log\log x}$ still tend to a Gaussian as $n$ runs through $(x,x+h]$? Proposition~\ref{prop:tkshort} does not reach those intervals, because it needs $y\le\sqrt h$ and the primes up to $y$ have to carry the variance $\log\log x$.

\section{Project: consecutive residues}

Chapter~\ref{ch:oc-bias} computes, under the generalised Riemann hypothesis and linear independence, the logarithmic density of the set of $x$ for which there are more primes up to $x$ in one residue class than in another. Lemke Oliver and Soundararajan observed a different bias: among consecutive primes, the pairs of residue classes are not equally likely, and the diagonal is light. That statement is not a theorem of this volume.

The singular series of a pair of residue classes $a,b$ modulo $q$ is positive and depends on $a-b$. An average of that series, if it were known in a range of moduli growing with $x$, would imply a bias for consecutive primes. The average the series has actually proved is the one-dimensional one: the race of two classes, integrated in the logarithmic variable $u=\log x$, under the two hypotheses named in Chapter~\ref{ch:oc-bias}. A numerical matrix of consecutive residues, however far it is computed, is a computation. It becomes the theorem only when the average of the singular series is proved, and that average is the open line.
